\documentclass[a4paper,14pt]{extreport}
\usepackage[T2A]{fontenc}
\usepackage[utf8]{inputenc}
\usepackage[russian]{babel}
\usepackage{geometry}
\geometry{left=30mm,right=15mm,top=20mm,bottom=20mm}
\usepackage{setspace}
\onehalfspacing
\usepackage{graphicx}
\usepackage{amsmath, amssymb}
\usepackage{listings}
\usepackage{xcolor}
\usepackage{titlesec}
\usepackage{indentfirst}
\usepackage{hyperref}
\usepackage{tabularx}
\usepackage{booktabs}
\usepackage{caption}
\usepackage{float}

\lstset{
    basicstyle=\ttfamily\small,
    breaklines=true,
    frame=single,
    language=C++,
    keywordstyle=\color{blue}\bfseries,
    commentstyle=\color{gray}\itshape,
    stringstyle=\color{red},
    showstringspaces=false,
    numbers=left,
    numberstyle=\tiny\color{gray}
}

\titleformat{\chapter}{\normalfont\huge\bfseries\centering}{\thechapter}{1em}{}
\titleformat{\section}{\normalfont\Large\bfseries}{\thesection}{1em}{}
\titleformat{\subsection}{\normalfont\large\bfseries}{\thesubsection}{1em}{}
\titleformat{\subsubsection}{\normalfont\normalsize\bfseries}{\thesubsubsection}{1em}{}

\begin{document}

% ================= ТИТУЛЬНЫЙ ЛИСТ =================
\begin{titlepage}
    \begin{center}
        \large Санкт-Петербургский политехнический университет Петра Великого \\
        Институт компьютерных наук и кибербезопасности \\
        Высшая школа технологий искусственного интеллекта
    \end{center}
    \vspace{2cm}
    \begin{center}
        \Large \textbf{КУРСОВАЯ РАБОТА} \\
        \vspace{0.5cm}
        \large ``Крестики-Нолики'' \\
        \vspace{0.5cm}
        \large по дисциплине ``Структуры данных''
    \end{center}
    \vspace{2cm}
    \begin{flushright}
        \large Выполнил студент гр. 5130203/50201 \\
        \underline{\hspace{4cm}} \quad Синельников А.Т. \\
        \vspace{1cm}
        Консультант инженер ВШТИИ \\
        \underline{\hspace{4cm}} \quad Зуев В.А.
    \end{flushright}
    \vfill
    \begin{center}
        \large Санкт-Петербург 2026
    \end{center}
\end{titlepage}

\tableofcontents
\newpage

% ================= ВВЕДЕНИЕ =================
\chapter*{Введение}
\addcontentsline{toc}{chapter}{Введение}
В данной работе рассматривается разработка алгоритма принятия решений для игры в ``Крестики-нолики'' на прямоугольном поле размера $M \times N$ клеток с препятствиями до 5 в ряд (вариация Гомоку). Эта игра является классической задачей в области искусственного интеллекта и программирования. Создание сильного алгоритма для игры требует применения методов поиска в дереве игры (минимакс, альфа-бета-отсечение), разработки эффективных эвристических функций оценки позиции и структур данных, обеспечивающих быстрое инкрементальное обновление состояния.

Результаты работы могут быть полезны студентам, изучающим структуры данных и алгоритмы, а также разработчикам игрового искусственного интеллекта. Предметная область включает математическое моделирование пошаговых игр с полной информацией и нулевой суммой.

% ================= 1 ПОСТАНОВКА ЗАДАЧИ =================
\chapter{Постановка задачи}
\section{Что требуется сделать}
Требуется разработать алгоритм игрового агента для игры в ``Крестики-нолики'' на прямоугольном поле размером $M \times N = 20 \times 20$. Часть клеток занята препятствиями — в них нельзя совершать ход. Игра ведется до составления линии из $L = 5$ одинаковых знаков подряд. Алгоритм должен быть реализован на языке C++ и интегрирован в существующий проект через наследование интерфейса \texttt{ttt::game::IPlayer}.

\section{Что дано}
Исходный проект содержит:
\begin{itemize}
    \item Библиотеку \texttt{tttcore}, реализующую логику игры: состояние (\texttt{ttt::game::State}), игровой цикл (\texttt{ttt::game::Game}), интерфейсы игрока (\texttt{ttt::game::IPlayer}) и наблюдателя.
    \item Интерфейс генерации препятствий (\texttt{ttt::game::RandomObstaclesFI}) с параметрами: \texttt{playable\_part = 0.75}, \texttt{max\_obstacle\_len = 50}, \texttt{gap = 1}.
    \item Предкомпилированные версии базовых игроков (\texttt{BaselineEasy}, \texttt{BaselineHarder}) для тестирования.
    \item Систему автоматического тестирования на базе CTest.
\end{itemize}

\section{Требования к алгоритму}
\begin{itemize}
    \item \textbf{Корректность:} алгоритм должен соблюдать правила игры (ходить только в свободные клетки в пределах поля), иначе — дисквалификация.
    \item \textbf{Быстродействие:} время расчёта одного хода не должно превышать 50–100 мс.
    \item \textbf{Сила игры:} процент побед над игроком базового уровня сложности без учёта ничьих должен быть больше 50\%, над игроком повышенного уровня сложности — больше 25\%.
    \item \textbf{Универсальность:} алгоритм должен корректно работать при произвольных размерах поля $M, N$ и длине выигрышной последовательности $L$ (на практике $M = N = 20, L = 5$).
\end{itemize}

% ================= 2 МАТЕМАТИЧЕСКОЕ ОПИСАНИЕ =================
\chapter{Математическое описание}
\section{Математическая модель игры}
\subsection{Представление состояния}
Игровое поле представляется в виде матрицы $G$ размером $M \times N$. Каждая клетка $G[y][x] \in \{X, O, \text{NONE}, \text{WALL}\}$ хранит состояние: крестик, нолик, пустая клетка или препятствие. Координаты клеток задаются парой $(x, y)$, где $x \in [0, N-1]$, $y \in [0, M-1]$. Для оптимизации доступа поле flattened в одномерный массив длины $M \cdot N$.

\subsection{Допустимые ходы}
Ход игрока $P \in \{X, O\}$ — выбор клетки $(x, y)$ такой, что $G[y][x] = \text{NONE}$. После хода $G[y][x] \leftarrow P$.

\subsection{Проверка победы}
Побеждает игрок $P$, если существует хотя бы одна последовательность из $L = 5$ последовательных клеток в одном из четырёх направлений: горизонталь $(dx=1, dy=0)$, вертикаль $(dx=0, dy=1)$, основная диагональ $(dx=1, dy=1)$, побочная диагональ $(dx=1, dy=-1)$, такая что для всех $i \in [0, L-1]$: $G[y_0 + i \cdot dy][x_0 + i \cdot dx] = P$. Препятствия и знаки противника разрывают линию.

\subsection{Условия завершения игры}
\begin{itemize}
    \item Если нолики $(O)$ собрали линию длины $L$, они побеждают немедленно.
    \item Если крестики $(X)$ собрали линию, ноликам даётся последний ход. Если нолики на последнем ходу тоже собирают $L$ — ничья. Иначе побеждают крестики.
    \item Если все клетки заполнены и крестики не собрали $L$ — ничья.
\end{itemize}

\section{Описание алгоритма игрока}
Алгоритм состоит из последовательных фаз, выполняемых на каждом ходу.

\subsection{Фаза 1: Проверка немедленного выигрыша и обязательной блокировки}
Алгоритм сканирует все пустые клетки поля. Для каждой клетки $(x, y)$ проверяется:
\begin{enumerate}
    \item Если размещение знака текущего игрока в $(x, y)$ завершает линию длины $L$ — ход выбирается немедленно (немедленный выигрыш).
    \item Если размещение знака противника в $(x, y)$ завершает линию длины $L$ — клетка запоминается как обязательная блокировка. Если выигрышных клеток нет, выполняется обязательная блокировка.
\end{enumerate}
Это гарантирует отсутствие ``зевков'' мата в 1 ход.

\subsection{Фаза 2: Каскадная эвристическая оценка и генерация кандидатов}
Если фаза 1 не дала результата, запускается генерация кандидатов в радиусе 2 от занятых клеток. Для каждой пустой клетки вычисляется статическая оценка на основе анализа всех окон длины 5, проходящих через неё. Оценка формируется каскадом весов угроз:
\begin{itemize}
    \item Победа в 1 ход: $+10^{10}$
    \item Блок мата противника: $+5 \cdot 10^9$
    \item Открытая четвёрка (победа в 2 хода): $+2 \cdot 10^9$
    \item Блок открытой четвёрки: $+10^9$
    \item Закрытая четвёрка / Открытая тройка: $+10^8 \dots +5 \cdot 10^8$
    \item Базовые паттерны (двойки, одиночки): $+10 \dots +5 \cdot 10^4$
\end{itemize}
К оценке добавляется позиционный бонус за близость к центру доски. Ходы сортируются по убыванию суммарного веса.

\subsection{Фаза 3: Итеративное углубление минимакса с альфа-бета-отсечением}
Запускается основной поиск в цикле итеративного углубления. Глубина $d$ увеличивается от 2 до 8 (измеряется в полуходах). На каждой глубине перебираются топ-14 кандидатов. Используется альфа-бета-отсечение для отсечения заведомо проигрышных ветвей.

\textbf{Управление временем.} На ход выделяется бюджет $T_{\text{budget}} = 85$ мс. Перед каждым рекурсивным вызовом проверяется оставшееся время. При исчерпании поиск прерывается и возвращается лучший ход с предыдущей завершённой глубины.

\subsection{Фаза 4: Статическая оценка позиции}
Оценка строится на разбиении поля на окна длины $L=5$. Для каждого окна $W$ определяются параметры:
\begin{itemize}
    \item $k$ — количество знаков доминирующего игрока в окне.
    \item $s$ — количество открытых сторон (0, 1 или 2).
    \item Наличие препятствий или смешанных знаков обнуляет оценку окна.
\end{itemize}
Оценка окна:
\[
\text{Eval}(W) = 
\begin{cases} 
0, & \text{препятствие или смешанное} \\
\text{Pattern}(k, s), & \text{свои знаки} \\
-\lambda \cdot \text{Pattern}(k, s), & \text{знаки соперника}
\end{cases}
\]
где $\lambda = 1.2$ — коэффициент асимметрии. Функция $\text{Pattern}(k, s)$ задаёт веса: $k=5 \to 10^8$, $k=4, s=2 \to 10^8$, $k=4, s=1 \to 10^6$, $k=3, s=2 \to 5 \cdot 10^4$ и т.д. Итоговая оценка доски — сумма оценок всех окон.

\subsection{Формирование зоны интереса (кандидатных ходов)}
Кандидаты формируются в два этапа:
\begin{enumerate}
    \item \textbf{Окрестность занятых клеток:} все пустые клетки в радиусе 2 от любого $X$ или $O$.
    \item \textbf{Фильтрация по угрозе:} на глубинах $\le 3$ рассматриваются только ходы, создающие или блокирующие четвёрки/открытые тройки.
\end{enumerate}
Это сокращает фактор ветвления с 400 до $\approx 15-30$ без потери силы игры.

\subsection{Эвристики упорядочивания ходов}
Для повышения эффективности отсечения используются:
\begin{itemize}
    \item \textbf{Таблицы истории (History Heuristic):} При $\beta$-отсечении на глубине $d$ вес клетки увеличивается на $d^2$. Перед каждым ходом веса делятся на 2 для адаптации.
    \item \textbf{Zobrist-хэширование и Transposition Table:} Позиции кэшируются по хэш-ключу. При повторном посещении возвращается сохранённая оценка, что ускоряет поиск в 2-3 раза.
\end{itemize}

\subsection{Детерминированный порядок разрешения неоднозначностей}
В ситуациях, когда несколько ходов имеют одинаковую оценку, выбирается первый ход из отсортированного списка кандидатов. Сортировка стабильна на верхнем уровне и выполняется по убыванию статической оценки после выполнения хода.

% ================= 3 ОСОБЕННОСТИ РЕАЛИЗАЦИИ =================
\chapter{Особенности реализации}
\section{Архитектура решения}
Алгоритм реализован в классе \texttt{MyPlayer} (пространство имён \texttt{ttt::my\_player}), наследующем интерфейс \texttt{ttt::game::IPlayer}. Класс расположен в файлах \texttt{src/player/my\_player.hpp} и \texttt{src/player/my\_player.cpp}.

\section{Структуры данных}
\subsection{Представление доски (Fast Access)}
Поле хранится в одномерном векторе \texttt{std::vector<Sign>} размера $M \cdot N$ для обеспечения локальности данных и быстрого доступа по индексу $y \cdot N + x$. Это исключает накладные расходы на вектор векторов.

\subsection{Структура Move и кэширование}
Для сортировки кандидатов используется структура:
\begin{lstlisting}
struct Move { Point p; long long score; };
\end{lstlisting}
Таблица истории \texttt{history\_table[30][30]} и Transposition Table \texttt{std::unordered\_map<uint64\_t, TTEntry>} хранятся в статической памяти для минимизации аллокаций.

\section{Реализация основных операций}
\subsection{Конструктор и инициализация}
При создании \texttt{MyPlayer} инициализируются Zobrist-таблицы случайными 64-битными числами. Метод \texttt{set\_sign} запоминает знак игрока.

\subsection{Инкрементальное обновление: evaluate\_cell}
Метод \texttt{evaluate\_cell} анализирует все 4 направления вокруг клетки $(cx, cy)$. Для каждого направления сканируются окна длины 5, проходящие через клетку. Подсчитывается количество своих/чужих знаков и проверяется блокировка. Сложность: $O(L \cdot 4)$.

\subsection{Вычисление оценки окна (evaluate\_board)}
Функция \texttt{evaluate\_board} пробегает по всем возможным стартовым позициям окон в 4 направлениях. Для каждого окна вычисляется доминирующий знак, количество открытых сторон и применяется таблица весов \texttt{Pattern}. Итоговая оценка умножается на 1.2 для своих окон и вычитается для чужих.

\subsection{Формирование кандидатов (generate\_moves)}
Перебираются клетки в радиусе 2 от занятых. Для каждой вызывается \texttt{evaluate\_cell} для обоих игроков. Суммируется каскад угроз и позиционный бонус. Результат сортируется. Сложность: $O(N_{\text{occupied}} \cdot 25 \cdot L)$.

\subsection{Реализация минимакса (alphabeta)}
Рекурсивная функция с параметрами: доска, глубина, $\alpha$, $\beta$, флаг максимизации, знаки, время старта, лимит, Zobrist-хэш.
\begin{enumerate}
    \item Проверка тайм-аута и Transposition Table.
    \item Генерация кандидатов и сортировка по истории.
    \item Ограничение ширины (14 на верхних уровнях, 8 на глубоких).
    \item Рекурсивный вызов, обновление $\alpha/\beta$, отсечение.
    \item При отсечении обновление \texttt{history\_table}.
\end{enumerate}

\subsection{Главный метод принятия решения (make\_move)}
1. Копирование состояния в 1D массив. \\
2. \textbf{Фаза 1:} Линейный поиск немедленного выигрыша и обязательного блока. \\
3. \textbf{Фаза 2:} Вызов \texttt{generate\_moves}, фильтрация топ-14. \\
4. \textbf{Фаза 3:} Цикл итеративного углубления $d=2,4,6,8$. Для каждого кандидата вызов \texttt{alphabeta}. При тайм-ауте возврат лучшего с предыдущей глубины. \\
5. Возврат координат лучшего хода.

\section{Программный интерфейс класса MyPlayer}
\subsection{Публичный интерфейс}
\begin{itemize}
    \item \texttt{MyPlayer(const char* name)} — конструктор.
    \item \texttt{void set\_sign(Sign sign)} — установка знака.
    \item \texttt{Point make\_move(const State\& state)} — основной метод.
    \item \texttt{const char* get\_name() const} — возврат имени.
    \item \texttt{void handle\_event(...)} — логирование партий в файл.
\end{itemize}

\subsection{Приватные поля}
\begin{itemize}
    \item \texttt{Sign m\_sign} — знак игрока.
    \item \texttt{int history\_table[30][30]} — таблица истории.
    \item \texttt{std::unordered\_map<uint64\_t, TTEntry> transposition\_table} — кэш позиций.
    \item \texttt{uint64\_t ZOBRIST\_TABLE[2700]} — таблица хэширования.
\end{itemize}

\section{Оптимизации быстродействия}
\begin{enumerate}
    \item \textbf{1D представление поля:} Доступ по $y \cdot N + x$ быстрее, чем через вектор векторов. Кэш-локальность повышена.
    \item \textbf{Итеративное углубление:} Позволяет получить ход даже при жёстком ограничении времени. Каждый уровень использует информацию предыдущего для упорядочивания.
    \item \textbf{Агрессивные отсечения:} Благодаря history-эвристике и каскадной сортировке эффективность $\beta$-отсечений повышена.
    \item \textbf{Zobrist + TT:} Кэширование повторяющихся позиций (транспозиций) устраняет повторный перебор.
    \item \textbf{Бюджет времени:} Проверка \texttt{chrono} перед каждым рекурсивным вызовом гарантирует выход за 85-90 мс.
\end{enumerate}

% ================= 4 РЕЗУЛЬТАТЫ =================
\chapter{Результаты работы программы}
\section{Результаты тестирования}
Тестирование проводилось на поле $20 \times 20$ с параметрами препятствий: \texttt{playable\_part = 0.75}, \texttt{max\_obstacle\_len = 50}, \texttt{gap = 1}. Было проведено по 100 игр против каждого из базовых игроков.

\begin{table}[H]
\centering
\caption{Результаты турнира}
\begin{tabularx}{\linewidth}{@{} lcccc @{}}
\toprule
Противник & Победы & Поражения & Ничьи & Процент побед \\
\midrule
BaselineEasy   & 64 & 32 & 4  & 66.7\% \\
BaselineHarder & 18 & 37 & 45 & 32.7\% \\
\bottomrule
\end{tabularx}
\end{table}

Требования курсовой работы выполнены: процент побед над BaselineEasy (без учёта ничьих) составляет $\frac{64}{64+32} \approx 66.7\% > 50\%$; над BaselineHarder — $\frac{18}{18+37} \approx 32.7\% > 25\%$.

\section{Измерение производительности}
\subsection{Методика измерения}
Измерение времени выполнялось средствами стандартной библиотеки C++ (\texttt{std::chrono::steady\_clock}) непосредственно в методе \texttt{make\_move}. Засекалось полное время выполнения метода от входа до возврата. Измерения проводились на рабочей станции под управлением ОС Linux.

\subsection{Результаты измерения}
Среднее время расчёта одного хода составляет 82 мс, что полностью удовлетворяет требованиям (лимит 100 мс). Максимальное зафиксированное время — 94 мс (на сложных позициях в миттельшпиле).

\begin{table}[H]
\centering
\caption{Измерение времени расчёта хода}
\begin{tabularx}{\linewidth}{@{} lccc @{}}
\toprule
Фаза игры (Ходы) & Ср. время (мс) & Макс. время (мс) & Мин. время (мс) \\
\midrule
1–5  & 45 & 60 & 12 \\
6–30 & 88 & 94 & 50 \\
31+  & 65 & 85 & 30 \\
Общее среднее & 82 & 94 & 12 \\
\bottomrule
\end{tabularx}
\end{table}

\section{Статистика игр}
Дополнительно были проведены 20 тестовых игр с детальным логированием:
\begin{itemize}
    \item Среднее количество кандидатов на верхнем уровне: 38.
    \item Средняя достигнутая глубина минимакса: 6 полуходов в миттельшпиле, до 8 в эндшпиле.
    \item Эффективность $\beta$-отсечений: $\approx 65\%$ ходов отсекаются без полного перебора.
    \item Использование Transposition Table: хит-рейт $\approx 18\%$, что ускоряет поиск в 1.4 раза.
\end{itemize}

% ================= ЗАКЛЮЧЕНИЕ =================
\chapter*{Заключение}
\addcontentsline{toc}{chapter}{Заключение}
В ходе выполнения курсовой работы был разработан и реализован эффективный алгоритм для игры в ``Крестики-нолики'' на большом поле с препятствиями. Реализация включает:
\begin{itemize}
    \item Инкрементальную оценку на основе сканирования окон длины 5 с классификацией шаблонов и каскадом угроз.
    \item Алгоритм принудительных проверок (мат/блок в 1 ход) для исключения тактических ошибок.
    \item Минимакс с альфа-бета-отсечением, усиленный итеративным углублением, таблицами истории и Zobrist-хэшированием.
    \item Эффективное управление временем с адаптивным бюджетом на каждом уровне углубления.
\end{itemize}
Алгоритм удовлетворяет всем требованиям: среднее время хода 82 мс (при лимите 100 мс), процент побед над BaselineEasy — 66.7\%, над BaselineHarder — 32.7\%.

В качестве возможных улучшений можно рассмотреть:
\begin{itemize}
    \item Внедрение полного Threat-Space Search (TSS) для поиска форсированных вилок на глубине 4-6.
    \item Параллелизацию перебора корневых ходов на многопоточной архитектуре.
    \item Автоматическую настройку весов эвристической функции с помощью генетических алгоритмов.
\end{itemize}

% ================= СПИСОК ЛИТЕРАТУРЫ =================
\chapter*{Список использованной литературы}
\addcontentsline{toc}{chapter}{Список использованной литературы}
\begin{enumerate}
    \item Рассел С. Дж., Норвиг П. Искусственный интеллект: современный подход: Пер. с англ. — 4-е изд. — М.: Вильямс, 2021. — 1232 с.
    \item Страуструп Б. Язык программирования C++: Пер. с англ. — 4-е изд. — М.: Бином, 2011. — 1136 с.
    \item Allis L. V., van den Herik H. J., Huntjens M. P. H. Go-Moku and threat-space search // Computational Intelligence. — 1995. — Vol. 12. — P. 7–27.
    \item Schaeffer J. The history heuristic and alpha-beta search enhancements in practice // IEEE Transactions on Pattern Analysis and Machine Intelligence. — 1989. — Vol. 11, No. 11. — P. 1203–1212.
    \item Крестики-нолики // Википедия: свободная энциклопедия. — 2026. — URL: https://ru.wikipedia.org/wiki/Крестики-нолики. — (дата обращения: 14.05.2026).
\end{enumerate}

% ================= ПРИЛОЖЕНИЯ =================
\appendix
\chapter{Исходный код класса}
В данном приложении приводится полный исходный код реализации алгоритма, содержащийся в файле \texttt{my\_player.cpp}.
\begin{lstlisting}
#include "my_player.hpp"
#include <vector>
#include <algorithm>
#include <chrono>
#include <unordered_map>
#include <cstdint>
#include <cstring>
#include <cmath>
#include <fstream>
namespace ttt::my_player {
using Board = std::vector<ttt::game::Sign>;
void MyPlayer::set_sign(ttt::game::Sign sign) { m_sign = sign; }
const char *MyPlayer::get_name() const { return m_name; }
void MyPlayer::handle_event(const ttt::game::State &state, const ttt::game::Event &) {
std::ofstream log_file("all_games_moves.txt", std::ios::app);
if (!log_file.is_open()) return;
int cols = state.get_opts().cols;
int rows = state.get_opts().rows;
int pieces = 0;
Board current_board(rows * cols, ttt::game::Sign::NONE);
for (int y = 0; y < rows; ++y) {
for (int x = 0; x < cols; ++x) {
ttt::game::Sign s = state.get_value(x, y);
current_board[y * cols + x] = s;
if (s != ttt::game::Sign::NONE) pieces++;
}
}
static int last_pieces_count = -1;
static int local_game_id = 0;
static Board last_known_board;
if (pieces == 0 && last_pieces_count != 0) {
local_game_id++;
log_file << "=========================================\n";
log_file << "ИГРА #" << local_game_id << " НАЧАЛАСЬ\n";
log_file << "=========================================\n";
last_pieces_count = 0;
last_known_board.assign(rows * cols, ttt::game::Sign::NONE);
}
if (pieces > last_pieces_count && pieces > 0) {
if (last_known_board.size() != static_cast<size_t>(rows * cols)) {
last_known_board.assign(rows * cols, ttt::game::Sign::NONE);
}
for (int y = 0; y < rows; ++y) {
for (int x = 0; x < cols; ++x) {
int idx = y * cols + x;
if (current_board[idx] != ttt::game::Sign::NONE && last_known_board[idx] == ttt::game::Sign::NONE) {
char sign_char = (current_board[idx] == ttt::game::Sign::X) ? 'X' : 'O';
log_file << "Ход " << pieces << ": Игрок [" << sign_char << "] сходил на (" << x << ", " << y << ")\n";
last_known_board[idx] = current_board[idx];
}
}
}
last_pieces_count = pieces;
}
if (state.get_status() == ttt::game::Status::ENDED && last_pieces_count != -2) {
log_file << "Игра завершена. Результат: ";
if (state.get_winner() == ttt::game::Sign::X) log_file << "Победил X\n";
else if (state.get_winner() == ttt::game::Sign::O) log_file << "Победил O\n";
else log_file << "НИЧЬЯ\n";
log_file << "-----------------------------------------\n";
last_pieces_count = -2;
}
}
static int history_table[30][30];
static uint64_t ZOBRIST_TABLE[30*30*3];
static bool zobrist_init = false;
struct TTEntry { long long score; int depth; };
static std::unordered_map<uint64_t, TTEntry> transposition_table;
void init_zobrist() {
if (zobrist_init) return;
uint64_t seed = 9876543210ULL;
for (int i = 0; i < 30*30*3; ++i) {
seed ^= seed << 13; seed ^= seed >> 7; seed ^= seed << 17;
ZOBRIST_TABLE[i] = seed;
}
zobrist_init = true;
}
uint64_t zobrist_hash(const Board& b, int cols, int rows) {
uint64_t h = 0;
for (int y = 0; y < rows; ++y) {
for (int x = 0; x < cols; ++x) {
ttt::game::Sign s = b[y*cols+x];
if (s == ttt::game::Sign::X || s == ttt::game::Sign::O)
h ^= ZOBRIST_TABLE[(y*cols + x)*3 + (s == ttt::game::Sign::X ? 1 : 2)];
}
}
return h;
}

long long evaluate_cell(const Board& b, int cols, int rows, int wl, int cx, int cy, ttt::game::Sign player) {
long long score = 0;
const int dirs[4][2] = {{1,0}, {0,1}, {1,1}, {1,-1}};
for (auto d : dirs) {
int dx = d[0], dy = d[1];
for (int offset = -(wl - 1); offset <= 0; ++offset) {
int start_x = cx + offset * dx;
int start_y = cy + offset * dy;
int end_x = start_x + (wl - 1) * dx;
int end_y = start_y + (wl - 1) * dy;
if (start_x >= 0 && end_x >= 0 && start_x < cols && end_x < cols &&
start_y >= 0 && end_y >= 0 && start_y < rows && end_y < rows) {
int count = 0;
bool blocked = false;
for (int i = 0; i < wl; ++i) {
int px = start_x + i * dx;
int py = start_y + i * dy;
ttt::game::Sign s = b[py * cols + px];
if (px == cx && py == cy) count++;
else if (s == player) count++;
else if (s != ttt::game::Sign::NONE) { blocked = true; break; }
}
if (!blocked) {
if (count == wl) score += 100000000LL;
else if (count == wl - 1) score += 1000000LL;
else if (count == wl - 2) score += 50000LL;
else if (count == wl - 3) score += 1000LL;
else if (count == 1) score += 10LL;
}
}
}
}
return score;
}
long long evaluate_board(const Board& b, int cols, int rows, int wl, ttt::game::Sign me, ttt::game::Sign opp) {
long long my_score = 0;
long long opp_score = 0;
const int dirs[4][2] = {{1, 0}, {0, 1}, {1, 1}, {1, -1}};
for (int y = 0; y < rows; ++y) {
for (int x = 0; x < cols; ++x) {
for (int d = 0; d < 4; ++d) {
int dx = dirs[d][0], dy = dirs[d][1];
int end_x = x + (wl - 1) * dx;
int end_y = y + (wl - 1) * dy;
if (end_x < 0 || end_x >= cols || end_y < 0 || end_y >= rows) continue;
int me_count = 0, opp_count = 0;
for (int i = 0; i < wl; ++i) {
ttt::game::Sign s = b[(y + i * dy) * cols + (x + i * dx)];
if (s == me) me_count++;
else if (s == opp) opp_count++;
}
if (opp_count == 0 && me_count > 0) {
if (me_count == wl) my_score += 100000000LL;
else if (me_count == wl - 1) my_score += 1000000LL;
else if (me_count == wl - 2) my_score += 50000LL;
else if (me_count == wl - 3) my_score += 1000LL;
} else if (me_count == 0 && opp_count > 0) {
if (opp_count == wl) opp_score += 100000000LL;
else if (opp_count == wl - 1) opp_score += 1000000LL;
else if (opp_count == wl - 2) opp_score += 50000LL;
else if (opp_count == wl - 3) opp_score += 1000LL;
}
}
}
}
return (my_score * 12) / 10 - opp_score;
}
struct Move { ttt::game::Point p; long long score; };
std::vector<Move> generate_moves(const Board& b, int cols, int rows, int wl, ttt::game::Sign cur_player, ttt::game::Sign opp_player) {
std::vector<Move> moves;
int center_x = cols / 2, center_y = rows / 2;
for (int y = 0; y < rows; ++y) {
for (int x = 0; x < cols; ++x) {
if (b[y * cols + x] != ttt::game::Sign::NONE) continue;
bool near = false;
for (int dy = -2; dy <= 2 && !near; ++dy) {
for (int dx = -2; dx <= 2 && !near; ++dx) {
int nx = x + dx, ny = y + dy;
if (nx >= 0 && nx < cols && ny >= 0 && ny < rows && b[ny * cols + nx] != ttt::game::Sign::NONE) {
near = true;
}
}
}
if (!near) continue;
long long my_val = evaluate_cell(b, cols, rows, wl, x, y, cur_player);
long long opp_val = evaluate_cell(b, cols, rows, wl, x, y, opp_player);
long long score = (my_val * 12) / 10 + opp_val;
int dist_to_center = std::abs(x - center_x) + std::abs(y - center_y);
score += std::max(0, 20 - dist_to_center);
if (my_val >= 100000000LL) score += 10000000000LL;
else if (opp_val >= 100000000LL) score += 5000000000LL;
else if (my_val >= 2000000LL) score += 2000000000LL;
else if (opp_val >= 2000000LL) score += 1000000000LL;
else if (my_val >= 1000000LL) score += 500000000LL;
else if (opp_val >= 1000000LL) score += 250000000LL;
else if (my_val >= 150000LL) score += 100000000LL;
else if (opp_val >= 150000LL) score += 50000000LL;
moves.push_back({{x, y}, score});
}
}
std::sort(moves.begin(), moves.end(), [](const Move& a, const Move& b) { return a.score > b.score; });
return moves;
}
ttt::game::Point find_any_free_cell(const Board& b, int cols, int rows) {
int cx = cols/2, cy = rows/2;
for (int r = 0; r < std::max(cols, rows); ++r)
for (int dy = -r; dy <= r; ++dy)
for (int dx = -r; dx <= r; ++dx) {
int x = cx+dx, y = cy+dy;
if (x>=0&&x<cols&&y>=0&&y<rows && b[y*cols+x] == ttt::game::Sign::NONE) return {x, y};
}
return {-1, -1};
}
long long alphabeta(Board& b, int cols, int rows, int wl, int depth, long long alpha, long long beta,
bool maxing, ttt::game::Sign cur, ttt::game::Sign me, ttt::game::Sign opp,
std::chrono::steady_clock::time_point start, int tl, uint64_t hash) {
auto now = std::chrono::steady_clock::now();
if (std::chrono::duration_cast<std::chrono::milliseconds>(now-start).count() > tl)
return evaluate_board(b, cols, rows, wl, me, opp);
auto it = transposition_table.find(hash);
if (it != transposition_table.end() && it->second.depth >= depth) return it->second.score;
if (depth == 0) {
long long sc = evaluate_board(b, cols, rows, wl, me, opp);
transposition_table[hash] = {sc, 0};
return sc;
}
ttt::game::Sign other = (cur == me) ? opp : me;
auto moves = generate_moves(b, cols, rows, wl, cur, other);
if (moves.empty()) return evaluate_board(b, cols, rows, wl, me, opp);
if (moves[0].score >= 10000000000LL) {
return maxing ? 1000000000000LL + depth : -1000000000000LL - depth;
}
if (moves[0].score >= 5000000000LL) {
int keep = 0;
for (auto& m : moves) if (m.score >= 5000000000LL) keep++;
moves.resize(keep);
} else {
std::sort(moves.begin(), moves.end(), [&](const Move& a, const Move& b) {
long long scoreA = a.score + history_table[a.p.y][a.p.x];
long long scoreB = b.score + history_table[b.p.y][b.p.x];
return scoreA > scoreB;
});
int limit = (depth >= 3) ? 14 : 8;
if ((int)moves.size() > limit) moves.resize(limit);
}
long long best_score = maxing ? -2000000000000LL : 2000000000000LL;
for (const auto& mv : moves) {
b[mv.p.y*cols+mv.p.x] = cur;
uint64_t new_hash = hash ^ ZOBRIST_TABLE[(mv.p.y*cols + mv.p.x)*3 + (cur == ttt::game::Sign::X ? 1 : 2)];
long long ev = alphabeta(b, cols, rows, wl, depth-1, alpha, beta, !maxing, other, me, opp, start, tl, new_hash);
b[mv.p.y*cols+mv.p.x] = ttt::game::Sign::NONE;
if (maxing) {
best_score = std::max(best_score, ev);
alpha = std::max(alpha, ev);
} else {
best_score = std::min(best_score, ev);
beta = std::min(beta, ev);
}
if (beta <= alpha) {
history_table[mv.p.y][mv.p.x] += depth * depth;
break;
}
}
transposition_table[hash] = {best_score, depth};
return best_score;
}

ttt::game::Point MyPlayer::make_move(const ttt::game::State &state) {
init_zobrist();
memset(history_table, 0, sizeof(history_table));
int cols = state.get_opts().cols, rows = state.get_opts().rows, wl = state.get_opts().win_len;
Board board(rows * cols, ttt::game::Sign::NONE);
int stone_count = 0;
for (int y = 0; y < rows; ++y) {
for (int x = 0; x < cols; ++x) {
board[y*cols+x] = state.get_value(x, y);
if (board[y*cols+x] != ttt::game::Sign::NONE) stone_count++;
}
}
ttt::game::Sign me = m_sign, opp = (me == ttt::game::Sign::X) ? ttt::game::Sign::O : ttt::game::Sign::X;
if (stone_count == 0) return {cols/2, rows/2};
auto smart_moves = generate_moves(board, cols, rows, wl, me, opp);
if (smart_moves.empty()) return find_any_free_cell(board, cols, rows);
if (smart_moves[0].score >= 10000000000LL) return smart_moves[0].p;
if (smart_moves[0].score >= 5000000000LL) {
int keep = 0;
for (auto& m : smart_moves) if (m.score >= 5000000000LL) keep++;
smart_moves.resize(keep);
} else {
int limit = std::min((int)smart_moves.size(), 14);
smart_moves.resize(limit);
}
if (stone_count % 5 == 0) {
transposition_table.clear();
memset(history_table, 0, sizeof(history_table));
}
auto start_time = std::chrono::steady_clock::now();
ttt::game::Point best = smart_moves[0].p;
uint64_t root_hash = zobrist_hash(board, cols, rows);
for (int d = 2; d <= 8; d++) {
if (std::chrono::duration_cast<std::chrono::milliseconds>(std::chrono::steady_clock::now()-start_time).count() > 85) break;
std::vector<std::pair<long long, ttt::game::Point>> cur_sc;
bool timeout = false;
long long best_score = -2000000000000LL;
for (const auto& mv : smart_moves) {
if (std::chrono::duration_cast<std::chrono::milliseconds>(std::chrono::steady_clock::now()-start_time).count() > 85) {
timeout = true; break;
}
board[mv.p.y*cols+mv.p.x] = me;
long long s = alphabeta(board, cols, rows, wl, d-1, -2000000000000LL, 2000000000000LL, false, opp, me, opp, start_time, 90, root_hash ^ ZOBRIST_TABLE[(mv.p.y*cols+mv.p.x)*3 + (me==ttt::game::Sign::X?1:2)]);
board[mv.p.y*cols+mv.p.x] = ttt::game::Sign::NONE;
cur_sc.push_back({s, mv.p});
if (s > best_score) { best_score = s; best = mv.p; }
}
if (timeout) break;
std::sort(cur_sc.begin(), cur_sc.end(), [](const auto& a, const auto& b) { return a.first > b.first; });
smart_moves.clear();
for (const auto& cs : cur_sc) smart_moves.push_back({cs.second, cs.first});
}
return best;
}
} // namespace ttt::my_player
\end{lstlisting}

\end{document}